\scanpage{163}
% L1–5
Ceci \inserted{(i.e. $\mathbf t$)} définit une involution sur
$E$ compatible avec l'augmentation, i.e.
une involution sur $A$ i.e. sur $\check{\mathbb V}(A)$,
laquelle laisse le diviseur $X$ invariant --
et y induit l'involution canonique.

% L6–19
Attention, si on se donne d'\uline{avance}
une involution sur $L$, et cherche les
div. quadr. sur $\check{\mathbb V}(L)$ qui \deleted{\unsure{soient stables|sont stables}} \inserted{la redonnent,}
\inserted{celles-ci sont stables} par l'involution, \deleted{ceci soit par \unsure{né.}}
mais cette condition n'est pas suffisante
-- en car. $2$, on peut prendre l'involution
identique, alors qu'il y aura des
diviseurs quadratiques (scindés, p.ex)
pour lesquels l'involution n'\uline{est} \uline{pas}
l'identité! Mais \deleted{\unsure{si|car}} $2$ inv., cette difficulté
disparaît -- cela se ramène à voir
que tt div. sur $\check{\mathbb V}(\deleted{\unsure{A}}L)\deleted{\unsure{E'|E}}$, stable
par $\lambda\to-\lambda$, est bien l'image inverse
d'un $\deleted{\unsure{\zeta}}\unsure{\zeta|\xi}\in L^2$ par $L\xrightarrow{u\mapsto u^{\otimes 2}} L^{\otimes 2}$.

% L20–24
(Il suffit $=$ que $2$ ne soit pas
div. de zéro \unsure{\ldots} \deleted{\unsure{?}} Il paraît que cette
dernière propriété soit vraie -- mais
pour l'appliquer aux $A$ affines \unsure{quelc.|quelq.}
il faut supposer que l'involution sur $A$
a un p$^{\mathrm t}$ fixe\ldots)
% archive sheet number: 162
